\documentclass[12pt]{article}
\usepackage[utf8]{inputenc}
\usepackage{amsmath}
\usepackage{geometry}
\geometry{a4paper, margin=2.5cm}

\title{Research on Image Classification Methods Based on Deep Learning}
\author{John Smith}
\date{\today}

\begin{document}

\maketitle

\begin{abstract}
This paper proposes a novel image classification method based on deep learning.
By introducing attention mechanisms and residual connections, our method achieves state-of-the-art performance on multiple benchmark datasets.
Experimental results demonstrate that the proposed method has good generalization ability and robustness.
\end{abstract}

\section{Introduction}

Image classification is a fundamental task in the field of computer vision.
With the development of deep learning technology, convolutional neural networks have made breakthrough progress in image classification tasks.

The main contributions of this paper include:
\begin{itemize}
    \item A novel network architecture is proposed
    \item An effective training strategy is designed
    \item Comprehensive evaluations are conducted on multiple datasets
\end{itemize}

\section{Related Work}

Deep learning has achieved great success in computer vision.
AlexNet first demonstrated the powerful capabilities of deep convolutional networks in the ImageNet competition.
Subsequently, network architectures such as VGGNet and ResNet further improved performance.

\section{Method}

\subsection{Network Architecture}

The proposed network architecture is shown in the figure.
The network consists of multiple convolutional layers and pooling layers, and finally outputs classification results through fully connected layers.

\subsection{Loss Function}

We use the cross-entropy loss function for training:
\begin{equation}
    L = -\sum_{i=1}^{N} y_i \log(p_i)
    \label{eq:loss}
\end{equation}
where $N$ is the number of categories, $y_i$ is the ground truth label, and $p_i$ is the predicted probability.

\section{Experiments}

\subsection{Datasets}

We conducted experiments on CIFAR-10, CIFAR-100, and ImageNet datasets.

\subsection{Experimental Results}

Experimental results show that our method achieves the best performance on all datasets.

\section{Conclusion}

This paper proposes a novel image classification method based on deep learning.
The effectiveness of the method has been validated through extensive experiments.
Future work will explore more network architecture designs and training techniques.

\end{document}

